% LaTeX report: FPL residual dataset generation pipeline
\documentclass[11pt]{article}
\usepackage[utf8]{inputenc}
\usepackage{amsmath,amssymb,bm}
\usepackage{geometry}
\usepackage{booktabs}
\usepackage{enumitem}
\geometry{margin=1in}

\title{FPL Residual Dataset Generation Pipeline}
\author{FPL-GNN Project}
\date{}

\begin{document}
\maketitle

\section{Purpose}

The dataset is designed to train a surrogate (e.g., GNN or MLP) to predict the \textbf{FPL residual}: the error of the linearized voltage--power model
\begin{equation}
  \mathbf{v}_{\mathrm{fpl}} = \mathbf{v}_{\mathrm{base}} + \mathbf{J} \begin{bmatrix} \Delta\mathbf{P} \\ \Delta\mathbf{Q} \end{bmatrix}
\end{equation}
relative to the true AC power-flow solution $\mathbf{v}_{\mathrm{true}}$. The target is
\begin{equation}
  \mathbf{r} = \mathbf{v}_{\mathrm{true}} - \mathbf{v}_{\mathrm{fpl}} \quad \text{(per-node voltage magnitude residual)}.
\end{equation}
All quantities are at \textbf{node level} (bus-phase); $\mathbf{J}$ has shape $(N \times 2N)$ with $N$ the number of nodes.

%==============================================================================
\section{Scenario and Time Resolution}
%==============================================================================

\subsection{Scenario-level draws}

For each scenario (one full day), four scalars are drawn \textbf{once} and held constant for the entire day:
\begin{itemize}[noitemsep]
  \item $P_{\mathrm{load}}$, $Q_{\mathrm{load}}$: total load (kW, kvar).
  \item $P_{\mathrm{der}}$, $Q_{\mathrm{der}}$: total DER level (kW, kvar).
\end{itemize}
Each is drawn from a Normal distribution with scenario-specific mean ($\mu_{P_{\mathrm{load}}}$, $\mu_{Q_{\mathrm{load}}}$, $\mu_{P_{\mathrm{der}}}$, $\mu_{Q_{\mathrm{der}}}}$) and a common standard deviation $\sigma_{\mathrm{shared}}$, then clipped to $\ge 0$. No per-timestep noise is applied when applying these to the circuit ($\sigma_{\mathrm{load}}=0$, $\sigma_{\mathrm{der}}=0$).

\subsection{Time resolution}

\begin{itemize}[noitemsep]
  \item \textbf{5-minute steps}: 288 steps per day, 12 steps per hour ($\mathtt{STEPS\_PER\_HOUR} = 12$).
  \item \textbf{Load} at step $t$: total load scaled by the load profile coefficient $m_L[t]$ (from \texttt{5minDayShape}).
  \item \textbf{DER} at step $t$: $P_{\mathrm{der}}(t) = P_{\mathrm{der}} \times m_{\mathrm{PV}}[t]$, $Q_{\mathrm{der}}(t) = Q_{\mathrm{der}} \times m_{\mathrm{PV}}[t]$, with $m_{\mathrm{PV}}$ from the irradiance/DER profile. The feeder maps this total DER to individual nodes (e.g., via PVSystem elements) so that per-node injections $\Delta P_i$, $\Delta Q_i$ sum to the desired totals; the same profile and mapping define both P and Q shares per node (e.g., $Q_i / P_i = Q_{\mathrm{der}}(t)/P_{\mathrm{der}}(t)$ when $P_{\mathrm{der}}(t) \neq 0$).
\end{itemize}

%==============================================================================
\section{Per-Hour Computation: Base and Jacobian}
%==============================================================================

\subsection{Load-only base (start of hour)}

At the \textbf{beginning of each hour} $h$:
\begin{enumerate}[noitemsep]
  \item Set the simulation time to the representative step for that hour: $\mathtt{step} = h \times \mathtt{STEPS\_PER\_HOUR}$.
  \item Apply \textbf{load only} (DER = 0): total $P_{\mathrm{load}}$, $Q_{\mathrm{load}}$ scaled by $m_L[\mathtt{step}]$; DER set to zero.
  \item Run one AC power-flow solve.
  \item If converged: record $\mathbf{v}_{\mathrm{base}}$ (voltage magnitude in pu at all $N$ nodes) and optional $\mathbf{x}_{\mathrm{base}}$ (nodal P/Q from load). If not converged, that hour is skipped (no samples for that hour).
\end{enumerate}

\subsection{Finite-difference Jacobian (same hour)}

Immediately after obtaining $\mathbf{v}_{\mathrm{base}}$ for that hour:
\begin{enumerate}[noitemsep]
  \item For each node $i = 0, \ldots, N-1$:
  \begin{itemize}[noitemsep]
    \item Add a new load at that node with $\mathtt{kW}=-1$, $\mathtt{kvar}=0$ (i.e., inject $+1$ kW); run AC solve; column $i$ of $\mathbf{J}$ is $(\mathbf{v}_{\mathrm{pert}} - \mathbf{v}_{\mathrm{base}}) / 1$.
    \item Add a new load at that node with $\mathtt{kW}=0$, $\mathtt{kvar}=-1$ (inject $+1$ kvar); run AC solve; column $N+i$ of $\mathbf{J}$ is $(\mathbf{v}_{\mathrm{pert}} - \mathbf{v}_{\mathrm{base}}) / 1$.
  \end{itemize}
  \item Result: $\mathbf{J}$ of shape $(N \times 2N)$ in units pu/kW and pu/kvar.
\end{enumerate}

\subsection{Fixed for the whole hour}

Both $\mathbf{v}_{\mathrm{base}}$ and $\mathbf{J}$ are computed \textbf{once at the start of the hour} and then \textbf{reused for all 12 five-minute steps} within that hour. No additional load-only or perturbation solves are performed for the intermediate timestamps.

%==============================================================================
\section{Per-Step Computation: True Voltage and Residual}
%==============================================================================

For each 5-minute step $t$ within the hour:
\begin{enumerate}[noitemsep]
  \item Apply \textbf{load + DER} at step $t$: load scaled by $m_L[t]$, DER by $m_{\mathrm{PV}}[t]$, with totals $(P_{\mathrm{load}}, Q_{\mathrm{load}})$ and $(P_{\mathrm{der}}(t), Q_{\mathrm{der}}(t))$ distributed by the feeder's device mapping.
  \item Run \textbf{one} AC power-flow solve $\to$ $\mathbf{v}_{\mathrm{true}}$ and, per node, the DER injection $\Delta P_i$, $\Delta Q_i$ (kW, kvar) at that step.
  \item Form $\boldsymbol{\delta}\mathbf{x} = [\Delta\mathbf{P}; \Delta\mathbf{Q}]$ (length $2N$).
  \item Linear FPL prediction (same $\mathbf{v}_{\mathrm{base}}$ and $\mathbf{J}$ for the whole hour):
  \begin{equation}
    \mathbf{v}_{\mathrm{fpl}} = \mathbf{v}_{\mathrm{base}} + \mathbf{J} \, \boldsymbol{\delta}\mathbf{x}.
  \end{equation}
  \item Residual (surrogate target):
  \begin{equation}
    \mathbf{r} = \mathbf{v}_{\mathrm{true}} - \mathbf{v}_{\mathrm{fpl}}.
  \end{equation}
  \item Reject the sample if any entry of $\mathbf{r}$ is non-finite or if $\mathbf{v}_{\mathrm{true}}$ lies outside allowed bounds (e.g., $[0.5, 1.5]$ pu); otherwise append one sample and $N$ node-level rows to the dataset.
\end{enumerate}

So for timestamps \textbf{between} two hour boundaries we solve \textbf{only} for load+DER to get $\mathbf{v}_{\mathrm{true}}$; $\mathbf{v}_{\mathrm{fpl}}$ is always computed from the \textbf{base voltage and Jacobian at the beginning of that hour}.

%==============================================================================
\section{Outputs and Storage}
%==============================================================================

\begin{itemize}[noitemsep]
  \item \textbf{gnn\_node\_index\_master.csv}: node name and node index for all $N$ nodes.
  \item \textbf{gnn\_sample\_meta.csv}: one row per kept sample: \texttt{sample\_id}, \texttt{scenario\_id}, \texttt{hour}, \texttt{step}.
  \item \textbf{gnn\_node\_features\_and\_targets.csv}: one row per (sample, node): \texttt{sample\_id}, \texttt{node}, \texttt{node\_idx}, \texttt{delta\_P\_kw}, \texttt{delta\_Q\_kvar}, \texttt{vmag\_true}, \texttt{vmag\_fpl}, \texttt{vmag\_resid}.
\end{itemize}
The surrogate is trained to predict \texttt{vmag\_resid} (or equivalently to refine $\mathbf{v}_{\mathrm{fpl}}$) from the injection features and graph structure.

%==============================================================================
\section{Solve Count (Current Implementation)}
%==============================================================================

Per hour: $1$ (load-only base) $+ 2N$ (finite-difference perturbations) $+ 12$ (load+DER per step) $= 13 + 2N$ power-flow solves.

Per scenario (24 hours): $24 \times (13 + 2N)$. For $N=95$: $24 \times 203 = 4872$ solves per scenario.

For $n_{\mathrm{scenarios}}$ scenarios: $n_{\mathrm{scenarios}} \times (312 + 48N)$ total solves (e.g., 500 scenarios $\Rightarrow$ 2,436,000 solves).

%==============================================================================
\section{Two Methods for Accelerating Dataset Generation}
%==============================================================================

\subsection{Reduced injection points and balanced three-phase assumption}

Instead of perturbing \textbf{every} node with $\pm 1$ kW and $\pm 1$ kvar (which yields $\mathbf{J}$ of size $N \times 2N$ and requires $2N$ perturbation solves per hour), one can restrict to a small set of \textbf{candidate buses} where DER is allowed to inject, and assume injections are \textbf{balanced three-phase}.

\begin{itemize}[noitemsep]
  \item \textbf{Perturbation}: 1 kW from a balanced three-phase DER at a bus $\Rightarrow$ 1 kW total, i.e., $1/3$ kW per phase (and similarly 1 kvar total $\Rightarrow$ $1/3$ kvar per phase if Q is included).
  \item \textbf{Two buses}: with two candidate buses, one measures the sensitivity of all node voltages to:
  \begin{itemize}[noitemsep]
    \item 1 kW ($0.33$ kW per phase) DER at bus 1,
    \item 1 kvar ($0.33$ kvar per phase) DER at bus 1,
    \item the same at bus 2.
  \end{itemize}
  \item This gives \textbf{4 sensitivity directions} $\Rightarrow$ $\mathbf{J}$ has \textbf{4 columns} (and $N$ rows): $\mathbf{J} \in \mathbb{R}^{N \times 4}$. The inputs to the linear model become the aggregated P and Q at the two buses only; $\boldsymbol{\delta}\mathbf{x}$ is 4-dimensional (e.g., $\Delta P_1$, $\Delta Q_1$, $\Delta P_2$, $\Delta Q_2$).
\end{itemize}

Per hour this requires $1$ (base) $+ 4$ (perturbations) $+ 12$ (load+DER) $= 17$ solves instead of $13 + 2N$, greatly reducing cost (e.g., for $N=95$, from 203 to 17 solves per hour). Dataset generation then stores the 4-dimensional injection vector per sample and the same $N$-dimensional residual target; the surrogate can be trained to predict the residual from these 4 inputs and the network structure.

\subsection{Analytical Jacobian instead of finite differences}

The Jacobian $\mathbf{J}$ can be computed \textbf{analytically} from the network admittance matrix and the FPL linearization (e.g., $\mathbf{Z}_{\mathrm{LL}}$ in per unit and the current voltage phasors), as implemented in the analytic--FD comparison module. That replaces the $2N$ (or 4) perturbation solves per hour with a single load-only solve (to get voltages and build $\mathbf{Y}$) and matrix operations to form $\mathbf{J}$. This further reduces the number of power-flow solves and is consistent with the same $\mathbf{v}_{\mathrm{base}}$ and $\mathbf{J}$ being used for the whole hour.

%==============================================================================
\section{Consistency with EMS and Accuracy Considerations}
%==============================================================================

Keeping $\mathbf{J}$ and $\mathbf{v}_{\mathrm{base}}$ \textbf{constant over the whole hour} is consistent with typical EMS implementations, where setpoints or sensitivity matrices are updated on a slower time scale (e.g., hourly) and applied over many shorter intervals. The dataset pipeline reflects this: one base and one Jacobian per hour, reused for all 12 steps in that hour.

However, this choice introduces a trade-off for the \textbf{power-flow surrogate}: the linear model $\mathbf{v}_{\mathrm{fpl}} = \mathbf{v}_{\mathrm{base}} + \mathbf{J}\boldsymbol{\delta}\mathbf{x}$ is \textbf{underspecified} over the hour in the sense that (i) load and DER profiles change at 5-minute resolution while the base and $\mathbf{J}$ are fixed at the hour start, and (ii) nonlinear and network effects (e.g., regulator taps, voltage-dependent load) are not captured by a single linearization. The residual $\mathbf{r} = \mathbf{v}_{\mathrm{true}} - \mathbf{v}_{\mathrm{fpl}}$ therefore carries both linearization error and the error due to using an hour-old base and Jacobian. Training the surrogate to predict this residual allows the combined FPL+surrogate model to improve accuracy while retaining the same hourly update structure and the option to accelerate data generation via reduced injection points or an analytical $\mathbf{J}$ as described above.

\end{document}
